\begin{abstract}
北半球冬季积雪变化涉及覆盖范围、雪水储量和覆盖边界，三者对变暖背景的响应可能不同。基于 1980--2022 年 ERA5 子集的 43 个冬季，固定分析域的积雪覆盖面积每十年减少 $0.575\times10^6$~km$^2$（95\% CI 为 $0.344$--$0.807\times10^6$~km$^2$ 的减幅），2 米气温每十年上升 $0.318\,^\circ$C。原序列平均雪水当量每十年减少 2.233~mm，但在模型中加入对应 2004 年积雪资料同化变更的水平项后，估计减幅缩小，区间跨零。覆盖面积和雪水当量去趋势后的相关接近零，标准化长期斜率差也未得到明确支持。以积雪覆盖率 0.50 定义的二维地理边界，其全分析域年度中位纬度每十年向北移动 $0.393^\circ$（95\% CI 为 $0.111^\circ$--$0.675^\circ$），阈值和经度支持敏感性检验的方向一致。这些结果表明，在该再分析子集中，覆盖收缩和边界北移的证据比雪水储量的长期下降更稳健。

\noindent\textbf{关键词}\quad 积雪覆盖率；雪水当量；冬季气温；阈值雪线；ERA5
\end{abstract}
